\documentclass[11pt]{article}
% Shared preamble for all daily exercises
% Update this file to change settings (padding, parindent, etc.) for all exercises

\usepackage[margin=0.75in]{geometry}
\usepackage{amsmath}
\usepackage{amssymb}

% Custom settings
\setlength{\parindent}{0pt}
\setlength{\parskip}{0.2cm}

\title{Daily Exercise}
\author{Dhruman Gupta}
\date{5th April}

\begin{document}

% Common header for daily exercises
% This file can be included in other LaTeX documents

\maketitle

\noindent
\textbf{Course:} CS-3330, Theory of Computation

\vspace{0.2cm}

\noindent
\textbf{Question:}\noindent 
The NFA square problem is: Does there exist a non-empty string $w$ such that $ww \in L(A)$?

Prove that this problem is NP-complete. 

\textbf{Answer:}

The statement is false as written. In fact, the NFA square problem is in $P$.

Let $A = (Q, \Sigma, \delta, I, F)$ be an NFA. We will show that there is a polynomial-time algorithm to decide whether there exists a non-empty string $w$ such that $ww \in L(A)$.

Observe that $ww \in L(A)$ iff there exists a state $q \in Q$ such that:

\begin{itemize}
    \item Starting from some initial state, $A$ can read $w$ and end in $q$.
    \item Starting from $q$, $A$ can read the same string $w$ and end in some final state.
\end{itemize}

So, the question is equivalent to asking whether there exists a state $q \in Q$ and a non-empty string $w$ such that
\[
w \in L(I, q) \cap L(q, F),
\]
where $L(I, q)$ is the set of strings taking some initial state to $q$, and $L(q, F)$ is the set of strings taking $q$ to some final state.

Now, fix a state $q$. Construct a product automaton whose states are pairs $(p, r) \in Q \times Q$. The first component simulates reading $w$ from an initial state, and the second component simulates reading the same $w$ from $q$.

The start states are all pairs $(i, q)$ for $i \in I$, and the accepting states are all pairs $(q, f)$ for $f \in F$.

For each symbol $a \in \Sigma$, put a transition
\[
(p, r) \xrightarrow{a} (p', r')
\]
whenever
\[
p' \in \delta(p, a) \text{ and } r' \in \delta(r, a).
\]

Then there exists a non-empty string $w$ such that
\[
i \xrightarrow{w} q \text{ and } q \xrightarrow{w} f
\]
for some $i \in I$ and $f \in F$ iff in this product automaton there is a path labelled by a non-empty string from some start state $(i, q)$ to some accepting state $(q, f)$.

Thus, for each $q \in Q$, we just need to test reachability in a graph of size $|Q|^2$. This is polynomial. Repeating for all $q \in Q$ is still polynomial.

Hence the NFA square problem is in $P$. Therefore, unless $P = NP$, it cannot be NP-complete.

\end{document}
